\clearpage
\begingroup
\section{Fixed-spin theta-channel super-Liouville sewing}
\label{app:partition-test}

The block in this partition-function calculation is exactly
$\mathbf F$ defined above and returned by the BPZ-native
double-Virasoro or superconformal $c$-recursion path. The
partition-function program does not append a sewing phase or a
second block conversion. The slots are $(\infty,1,0)$ and the edge
order is $(1,2,3)$. The form signs $\eta,\eta'$ and tube signs
$\eta_e$ are different labels. Each chosen set of tube signs
represents one fixed marked spin structure; no spin structures are
averaged.

Write $P_e=ip_e$ with $p_e>0$. The central charge, weights and
Ramond ground parameter are
\begin{equation}
 Q=b+b^{-1},\qquad c=\frac32+3Q^2,\qquad
 h_{\NS}(p)=\frac{Q^2}{8}+\frac{p^2}{2},\quad
 h_{\R}(p)=h_{\NS}(p)+\frac1{16},\quad
 \beta=\frac{ip}{\sqrt2}.
 \label{eq:partition-spectrum}
\end{equation}
The descendant blocks omit the primary factors $q_e^{h_e}$.
The continuum measure is $dp_e/\pi$ and the NS and physical Ramond
two-point coefficients are both $\pi\delta(p-p')$.

\subsection{Ordered forms and BPZ sewing}

The first-slot BPZ inversion is $z\mapsto1/z$ with lift $-i/z$;
each edge is sewn by $z\mapsto q_e/z$ with lift
$-i\sqrt{q_e}/z$. The linear graded transpose and ground metrics
are \eqref{eq:app-adjoint}--\eqref{eq:app-ground-metrics}.
In particular, the physical Ramond ground metric is
$\operatorname{diag}(1,-1)$ and the auxiliary one is
$\operatorname{diag}(1,-i)$.

Keep the ordered ground values, including
$\rho_1^{(\eta)}(\phi,w^-,w^+)=i\eta$. The square-root branch at
the origin is $i\sqrt z\sqrt{1-z}$, continued along the upper lip.
For an even NS primary the $m=n=0$ supercurrent identity reads
\begin{align}
 \rho_f^{(\eta)}(G_{-1/2}\phi,w_2^\alpha,w_3^\gamma)
 ={}&i\rho_f^{(\eta)}(\phi,G_0w_2^\alpha,w_3^\gamma)\nonumber\\
 &-(-1)^{f+|\alpha|}
       \rho_f^{(\eta)}(\phi,w_2^\alpha,G_0w_3^\gamma).
 \label{eq:partition-ordered-ward}
\end{align}
All other descendants use the four identities in
Appendix~\ref{app:wards}.

For an auxiliary-first tensor product the enlarged ordered
three-point form contains exactly
\begin{equation}
 (-1)^{\epsilon'_1\epsilon_1+\epsilon_2\epsilon'_3}
 \rho_{\F}(x'_1,x'_2,x'_3)\rho_f^{(\eta)}(x_1,x_2,x_3).
 \label{eq:partition-enlarged-vertex}
\end{equation}
The first exponent is graded BPZ dualization at infinity; the
second moves the middle physical state past the third auxiliary
state. Every tube uses the inverse of the full graded tensor Gram.
There is no separate $(-i)^{\epsilon_e}$ factor on a tube or
$i^{\epsilon_{v,1}}$ factor at an infinity slot.
The partition-function driver evaluates the returned coefficients
at the fixed $\eta_e$ and the chosen $q_e$, then applies only the
structure constants and primary factors displayed below.

\subsection{The three-point coefficients}


For an all-NS sphere, $C_a(p_1,p_2,p_3)$ is the coefficient of the
normalized holomorphic and antiholomorphic forms $\rho_a$.
Thus $C_0$ is the primary three-point coefficient, and $C_1$
is the coefficient with a $G_{-1/2}$ and an antiholomorphic
$G_{-1/2}$ on the middle primary, in the ordered convention of
the paper. The operator-moving signs remain in the sewing.
Here $C_0$ is real and $C_1$ is purely imaginary on the continuum.

For NS--R--R we denote the reduced SCFT coefficient in the sewing
formula below by $C_{f,\eta}$. The four normalized forms
$\rho_f^{(\eta)}$ do not give four independent Liouville couplings.
With the fixed Ramond ground-state factors kept in the sewing,
\begin{equation}
 C_{1,\eta}(p_1,p_2,p_3)=C_{0,\eta}(p_1,p_2,p_3).
 \label{eq:partition-Cf}
\end{equation}
The same sign-labeled coefficient multiplies the even and odd chiral
vertices in \href{https://arxiv.org/pdf/0810.1203}{Hadasz--Jask\'olski--Suchanek,
Eq.~(4.13)}; changing $f$ does not exchange $\eta$ with $-\eta$.
The current long draft states $C_{1,\eta}=C_{0,-\eta}$, which is
inconsistent with the coefficient assignment to its displayed
$\rho_1^{(\eta)}$ ground values and the fixed-spin decomposition used
here. That sentence needs correction
when the draft is next edited; the draft has not been changed here.
The two physical equal-family ground three-point values are
$C_{0,+}+C_{0,-}$ and $C_{0,+}-C_{0,-}$.
The label $f$ still distinguishes the normalized chiral forms;
their relative phases are not absorbed into $C_{f,\eta}$.

The explicit coefficients use the standard functions
\begin{equation}
 \Upsilon_{\NS}(x)=\Upsilon_b(x/2)\Upsilon_b((x+Q)/2),\qquad
 \Upsilon_{\R}(x)=\Upsilon_b((x+b)/2)\Upsilon_b((x+b^{-1})/2).
 \label{eq:partition-upsilon}
\end{equation}
These are the only additional special functions needed here.
For $0<\operatorname{Re}x<Q$ their normalization is fixed by
\begin{equation}
 \log\Upsilon_b(x)=\int_0^\infty\frac{dt}{t}
 \left[\left(\frac Q2-x\right)^2e^{-t}
 -\frac{\sinh^2((Q/2-x)t/2)}
 {\sinh(bt/2)\sinh(t/(2b))}\right],\qquad
 \Upsilon'_{\NS}(0)=\frac{\Upsilon_b(b)}2.
 \label{eq:partition-upsilon-definition}
\end{equation}
Analytic continuation outside this strip is fixed by
$\Upsilon_b(Q-x)=\Upsilon_b(x)$ and
\begin{equation}
 \Upsilon_b(x+b)=b^{1-2bx}
 \frac{\Gamma(bx)}{\Gamma(1-bx)}\Upsilon_b(x),\qquad
 \Upsilon_b(x+b^{-1})=b^{2x/b-1}
 \frac{\Gamma(x/b)}{\Gamma(1-x/b)}\Upsilon_b(x).
\end{equation}

With the common momentum-independent cosmological prefactor
fixed to one, as in the numerical test, the all-NS coefficients are
\begin{align}
 C_0(p_1,p_2,p_3)
 &=\frac{\Upsilon'_{\NS}(0)}{b^3}
 \frac{\displaystyle\prod_{j=1}^3
       \sqrt{\Upsilon_{\NS}(2ip_j)\Upsilon_{\NS}(-2ip_j)}}
 {\displaystyle\Upsilon_{\NS}(Q/2+i\sum_jp_j)
       \prod_{j=1}^3\Upsilon_{\NS}(Q/2+i\sum_kp_k-2ip_j)},\nonumber\\
 C_1(p_1,p_2,p_3)
 &=\frac{2i\Upsilon'_{\NS}(0)}{b^3}
 \frac{\displaystyle\prod_{j=1}^3
       \sqrt{\Upsilon_{\NS}(2ip_j)\Upsilon_{\NS}(-2ip_j)}}
 {\displaystyle\Upsilon_{\R}(Q/2+i\sum_jp_j)
       \prod_{j=1}^3\Upsilon_{\R}(Q/2+i\sum_kp_k-2ip_j)}.
 \label{eq:partition-Ca}
\end{align}
All momentum sums in \eqref{eq:partition-Ca} run from 1 to 3.
For NS--R--R, with $p_1$ on the NS edge, the two independent
coefficients are
\begin{align}
 C_{0,+}(p_1,p_2,p_3)
 &=\frac{\Upsilon'_{\NS}(0)}{b^3}
 \frac{\displaystyle\sqrt{\prod_{\sigma=\pm1}
     \Upsilon_{\NS}(2i\sigma p_1)
     \Upsilon_{\R}(2i\sigma p_2)\Upsilon_{\R}(2i\sigma p_3)}}
 {\displaystyle\prod_{\sigma=\pm1}
     \Upsilon_{\R}(Q/2+i(p_2+p_3+\sigma p_1))
     \Upsilon_{\NS}(Q/2+i(p_1+\sigma(p_2-p_3)))},\nonumber\\
 C_{0,-}(p_1,p_2,p_3)
 &=\frac{\Upsilon'_{\NS}(0)}{b^3}
 \frac{\displaystyle\sqrt{\prod_{\sigma=\pm1}
     \Upsilon_{\NS}(2i\sigma p_1)
     \Upsilon_{\R}(2i\sigma p_2)\Upsilon_{\R}(2i\sigma p_3)}}
 {\displaystyle\prod_{\sigma=\pm1}
     \Upsilon_{\NS}(Q/2+i(p_2+p_3+\sigma p_1))
     \Upsilon_{\R}(Q/2+i(p_1+\sigma(p_2-p_3)))}.
 \label{eq:partition-Cfeta}
\end{align}
All numerator square roots are positive for
$p_j>0$. The factor $b^{-3}$ includes the
normalized external legs: their gamma-function moduli are $1/b$
per NS leg and one per R leg, and the NS--R--R expression has
the additional relative factor $b^{-2}$. No separate leg
normalizations or three-point symbols are needed.
These formulas use the equal NS/R two-point convention of
\href{https://arxiv.org/pdf/hep-th/9607120}{Poghossian,
Eqs.~(14)--(15), (51), and (56)}.
The standalone Type 0B amplitude code uses half of each displayed
NS--R--R coefficient and a Ramond continuum two-point metric half as
large. Rescaling its Ramond field by $\sqrt2$ gives the equal-metric
convention here: each NS--R--R vertex gains a factor two and each
inverse Ramond continuum metric loses a factor two. On the theta
channel's two Ramond internal edges these factors cancel. This field
normalization changes no normalized chiral block.

\subsection{Fixed-spin decompositions}

For either theta chart, set the tube signs to the chosen fixed spin
before evaluating a block. The all-NS contribution is
\begin{equation}
 \left.Z_{\mathrm{SL}}\right|_{\NS\NS\NS,\eta_e}
 =\int_{(0,\infty)^3}\prod_{e=1}^3\frac{dp_e}{\pi}|q_e^{h_e}|^2
 \sum_{a=0}^1(-1)^a C_a^2
 \mathbf F^{(a)}(q;\eta_e)\overline{\mathbf F^{(a)}(q;\eta_e)}.
 \label{eq:partition-allns}
\end{equation}
The $i$ in $C_1$ and the $(-1)^a$ are both retained.
For the NS--R--R chart the physical Ramond ground pairing gives
\begin{align}
 \left.Z_{\mathrm{SL}}\right|_{\NS\R\R,\eta_e}
 ={}&\frac14\int\prod_{e=1}^3\frac{dp_e}{\pi}|q_e^{h_e}|^2
 \sum_{f=0}^1\sum_{\eta=\pm1}C_{f,\eta}^2
 \mathbf F_f^{(\eta,\eta)}(q;\eta_e)
 \overline{\mathbf F_f^{(\eta,\eta)}(q;\eta_e)}.
 \label{eq:partition-nsrr-full}
\end{align}
This sums normalized three-point forms at one fixed spin, not
surface spin structures. On the two even Ramond cycles tested,
$\mathbf F_1^{(\eta,\eta)}=-\mathbf F_0^{(\eta,\eta)}$ at the
\emph{same} $\eta$. With $C_{1,\eta}=C_{0,\eta}$,
\eqref{eq:partition-nsrr-full} reduces to
\begin{equation}
 \left.Z_{\mathrm{SL}}\right|_{\NS\R\R,\eta_e}
 =\frac12\int\prod_{e=1}^3\frac{dp_e}{\pi}|q_e^{h_e}|^2
 \sum_{\eta=\pm1}C_{0,\eta}^2
 \mathbf F_0^{(\eta,\eta)}(q;\eta_e)
 \overline{\mathbf F_0^{(\eta,\eta)}(q;\eta_e)}.
 \label{eq:partition-nsrr}
\end{equation}
The factor $1/2$ is the physical Ramond ground-pairing
normalization. It is fixed, not fitted to the cross-channel ratio.
In the NS identity limit $p_1\to iQ/2$ with $p_2=p_3$,
\begin{equation}
 \frac{C_{f,+}(p_1,p_2,p_2)}{C_0(p_1,p_2,p_2)}\to1,
 \qquad
 \frac{C_{f,-}(p_1,p_2,p_2)}{C_0(p_1,p_2,p_2)}\to0
 \quad(f=0,1).
 \label{eq:partition-identity-residue}
\end{equation}
The identity residue and the direct special-function evaluation
fix the relative NS/R normalization.

\subsection{Directed checks and scope}

The BPZ-native theta test at 40 digits compares direct physical
PBW sewing, the double-Virasoro recovery, and the coefficientwise
enlarged convolution through total level three. It covers both
intrinsic NS primary parities and the ordinary and inserted
Ramond-loop sectors. The local NS--R--R $G$ and auxiliary $\psi$
residue identities are separately tested in
\path{C++/tests/bpz_one_over_z_ward.cpp}; the tensor pairing
and enlarged vertex are tested in
\path{C++/tests/bpz_one_over_z_tensor.cpp}.
The separate theta PBW oracle is checked by
\path{C++/tools/check_theta_pbw_oracle.py}. It is an independent
theta implementation of the SCblock notes, not a block from
Yutai's sphere or torus calculation.

The fixed-spin partition comparison was rerun with native BPZ
blocks at 40 digits on ten matched genus-two surfaces. Each
surface uses two individually transported spin assignments.
The NS--R--R source has total descendant level five and
$7^3$ momentum nodes; the all-NS target has level eight and
$10^3$ nodes. The twenty separately evaluated source/target
ratios lie between $0.9998405584$ and $1.0000233972$; their
largest deviation from one is $1.59442\times10^{-4}$.
The run took $732.6$ seconds. Finite quadrature, geometry and
structure-constant accuracy contribute to this residual.
No spin sum and no coefficient-dependent sewing conversion
are used in the partition driver.

Yutai's compared data are actual sphere four-point and torus
two-point blocks. The sphere comparison uses common ordered
three-point forms and checks 864 coefficients through internal
level two, with maximum absolute discrepancy
$3.89\times10^{-16}$; the frozen physical-bank comparison
has maximum discrepancy $7.2\times10^{-9}$ on its stated
supported nodes. The torus comparison checks 576 open-R-edge
coefficients through NS twice-level two and R level one,
including the supported external descendants. It agrees to
$2.78\times10^{-16}$ after the BPZ sewn-edge frame phase
$i^{\epsilon_{\NS}-\epsilon_{\R}}$. This phase compares the
two implementations' local frames; it is not an extra factor
in the block defined here. These are block comparisons, not
recalculations of complete integrated Type 0B amplitudes.

The current directed commands are
\begin{verbatim}
make -C C++ check-partition
python3 C++/tools/check_theta_pbw_oracle.py
python3 C++/tools/check_yutai_sphere_blocks.py
python3 C++/tools/check_yutai_torus_open_bpz.py
\end{verbatim}
The implementation is in \path{C++/include/ramond/pipeline.hpp},
\path{C++/include/ramond/direct_pbw.hpp},
\path{C++/include/scblocks/ns_recursion.hpp}, and
\path{C++/drivers/partition_main.cpp}.
\endgroup
